\chapter{Network Cards and Kernel Bypass}\label{ch:nw:network-cards-and-kernel-bypass}

Between the wire and a UDP socket's \texttt{recv} a datagram is copied by the card into the host's memory, announced by an
interrupt, picked up by the kernel's networking code in a software interrupt, matched to a socket, queued, and copied again into
the buffer the program handed to \texttt{recv}; the program, meanwhile, was asleep and must be woken. On the laptop this book is
written on, the median datagram spends under a microsecond getting into the kernel over loopback and about fifteen more
microseconds getting out of it to a sleeping reader. None of that work decides anything a trading program cares about.

This chapter follows a packet through the network card and the kernel, then removes the kernel from the path. One Quant
Book~13 (chapter~13) tuned the host so that busy polling and kernel bypass can do their job; this chapter looks at the same
machinery from the card's side: \omterm{def:nw:network-cards-and-kernel-bypass:ring}{descriptor rings}, interrupts and their coalescing, the steering of flows to queues, the stacks
that run in user space, and the timestamps the card itself can take.

\section{The path of a packet through the kernel}

A network card and its driver share memory. The card writes arriving packets into buffers that the driver has given it,
and says which ones it has filled; the driver takes them, hands the buffers back, and passes the packets up.

\begin{definition}[Descriptor ring]\label{def:nw:network-cards-and-kernel-bypass:ring}
A \emph{descriptor ring}\index{descriptor ring} is a circular array of descriptors, each pointing at a packet buffer in host memory
and carrying an ownership flag. On the receive side the card fills the next descriptor it owns with an arriving packet (by direct
memory access), marks it as the host's, and advances its head; the host processes the descriptors it owns in order and hands them
back to the card. When the next descriptor still belongs to the host, the card has nowhere to put the packet and drops it.
\end{definition}

\begin{omfigure}
\begin{tikzpicture}[font=\footnotesize,b/.style={draw,rounded corners=2pt,minimum height=0.6cm,align=center,inner sep=3pt},>=Stealth]
\node[b,fill=gray!10] (w) at (0,0) {wire};
\node[b,fill=omDef!18] (n) at (2.0,0) {card};
\node[b,fill=omDef!10] (r) at (4.4,0) {descriptor\\ ring (DMA)};
\node[b,fill=omThm!12] (i) at (6.9,0.9) {interrupt,\\ softirq};
\node[b,fill=omThm!12] (s) at (9.2,0.9) {socket\\ queue};
\node[b,fill=omThm!12] (c) at (11.4,0.9) {copy, wake,\\ \texttt{recv}};
\node[b,fill=omMeth!20] (u) at (8.1,-1.1) {user-space poll of the ring (bypass)};
\draw[->] (w) -- (n);
\draw[->] (n) -- (r);
\draw[->,omThm,thick] (r.north east) -- (i.west);
\draw[->,omThm,thick] (i) -- (s);
\draw[->,omThm,thick] (s) -- (c);
\draw[->,omMeth,thick] (r.south east) -- (u.west);
\node[font=\scriptsize,omThm] at (9.2,1.75) {kernel path};
\node[font=\scriptsize,omMeth] at (8.1,-1.75) {bypass path: the program owns the ring};
\end{tikzpicture}

\omcaption{Two ways out of a \omterm{def:nw:network-cards-and-kernel-bypass:ring}{descriptor ring}. The kernel path (top) takes an interrupt, processes the packet in a software
interrupt, queues it on a socket and copies it to a program that must be woken. A kernel-bypass stack (bottom) maps the ring into
the program, which polls it and reads the packet where the card wrote it.}\label{fig:nw:network-cards-and-kernel-bypass:path}
\end{omfigure}

On Linux the kernel side of this is NAPI, the networking stack's event mechanism: the card raises an interrupt, the kernel
schedules the ring's processing in a software interrupt, and the processing runs until the ring is empty or a budget is spent,
with further interrupts from that ring disabled meanwhile. Under load the kernel is thus already polling; what it does not do by
default is spare the program the socket queue, the copy and the wake-up.

\begin{omfigure}
\begin{tikzpicture}[font=\scriptsize,>=Stealth]
\foreach \k in {0,...,11} {
  \pgfmathsetmacro{\a}{90-30*\k}
  \pgfmathtruncatemacro{\host}{(\k>=2 && \k<=6) ? 1 : 0}
  \ifnum\host=1 \def\c{omThm!30}\else\def\c{omDef!15}\fi
  \draw[fill=\c] (\a+15:1.3) arc (\a+15:\a-15:1.3) -- (\a-15:2.2) arc (\a-15:\a+15:2.2) -- cycle;
}
\draw[->,very thick,omMeth] (0,0) -- (30:1.2);
\node[font=\scriptsize,omMeth] at (30:0.45) [above left] {tail};
\draw[->,very thick,omDef] (0,0) -- (-120:1.2);
\node[font=\scriptsize,omDef] at (-120:0.5) [left] {head};
\node[anchor=west,align=left] at (2.8,0.9) {\tikz\fill[omThm!30] (0,0) rectangle (0.3,0.2);\ owned by the host: filled by the card,\\
  \hphantom{\tikz\fill (0,0) rectangle (0.3,0.2);\ }waiting to be processed or handed back};
\node[anchor=west,align=left] at (2.8,-0.3) {\tikz\fill[omDef!15] (0,0) rectangle (0.3,0.2);\ owned by the card: empty buffers};
\node[anchor=west,align=left] at (2.8,-1.2) {the card fills at the head and advances it clockwise;\\ the host reads at the tail;
  when the head reaches\\ a host-owned descriptor, the next packet is dropped};
\end{tikzpicture}

\omcaption{A receive ring of twelve descriptors. Five are the host's (filled and not yet handed back), seven are the card's. The card
writes at the head, the host reads at the tail; the ring is full when the head meets a descriptor the host still owns, whether the
host has read it or merely not yet handed it back.}\label{fig:nw:network-cards-and-kernel-bypass:ring}
\end{omfigure}

\begin{proposition}[When a ring overflows]\label{prop:nw:network-cards-and-kernel-bypass:full}
Let a ring of $N$ descriptors receive packets at rate $\lambda$ while the host processes none for a time $T$, starting from a ring in
which the host holds $H$ descriptors. The ring drops $\bigl(\lambda T - (N - H)\bigr)^{+}$ packets. A host that hands descriptors
back in batches of $b$ holds up to $b - 1$ processed descriptors at any time, so it must size the ring for the longest pause plus
$b - 1$.
\end{proposition}

\begin{proof}
The card has $N - H$ descriptors to fill; each arrival consumes one and none come back during the pause. Descriptors processed but
not yet handed back remain the host's until the batch completes.
\end{proof}

\section{Interrupts, coalescing and receive-side scaling}

An interrupt costs the host a switch into the kernel, the handler, and, for a sleeping reader, the scheduler's wake-up: microseconds,
paid per interrupt. Cards therefore wait before interrupting.

\begin{definition}[Interrupt coalescing]\label{def:nw:network-cards-and-kernel-bypass:coalesce}
\emph{Interrupt coalescing}\index{interrupt coalescing} is a card's practice of delaying the receive interrupt after a packet
arrives, until a given time has passed or a given number of packets are waiting, so that one interrupt serves several packets.
On Linux the two limits are set with \texttt{ethtool -C} as \texttt{rx-usecs} and \texttt{rx-frames}; zero microseconds and one frame
disable coalescing.
\end{definition}

\begin{proposition}[What coalescing buys and costs]\label{prop:nw:network-cards-and-kernel-bypass:coalesce}
Let packets arrive as a Poisson process of rate $\lambda$, let the card interrupt $u$ after the first waiting packet (the frame limit
never binding), and let each interrupt cost $c_{\mathrm{irq}}$ of CPU and each packet $c_{\mathrm{pkt}}$. Each interrupt serves on
average $1 + \lambda u$ packets, the host spends a fraction
\[
  \frac{\lambda\, c_{\mathrm{irq}}}{1 + \lambda u} + \lambda\, c_{\mathrm{pkt}}
\]
of one core, and the first packet of each batch waits $u$ for the interrupt. Without coalescing ($u = 0$) the host saturates one core
at $\lambda = 1/(c_{\mathrm{irq}} + c_{\mathrm{pkt}})$.
\end{proposition}

\begin{proof}
A batch opens with a packet and lasts $u$; by the memorylessness of the Poisson process, $\lambda u$ more packets arrive on
average. Interrupts therefore come at rate $\lambda/(1+\lambda u)$, each costing $c_{\mathrm{irq}}$, and every packet costs
$c_{\mathrm{pkt}}$.
\end{proof}

\begin{definition}[Receive-side scaling, flow steering]\label{def:nw:network-cards-and-kernel-bypass:rss}
\emph{Receive-side scaling}\index{receive-side scaling} (RSS) is a multi-queue card's distribution of arriving packets over several
receive rings, by a hash of their addresses and ports, so that different cores can process different flows. \emph{Flow
steering}\index{flow steering} is the card's placing of chosen flows (by exact match on addresses and ports) on a chosen ring,
overriding the hash.
\end{definition}

A trading host uses both against RSS's intent. RSS spreads load evenly, which is what a web server wants; a feed handler wants every
packet of line A on the one ring its pinned core polls, and nothing else there. \omterm{def:nw:network-cards-and-kernel-bypass:rss}{Flow steering} puts the venue's groups on the rings of
the cores that handle them and the rest of the machine's traffic elsewhere, and the core that polls a ring should sit on the socket
the card is attached to (One Quant Book~13, chapter~4).

\section{Polling drivers and user-space stacks}

\begin{definition}[Poll-mode driver, user-space network stack]\label{def:nw:network-cards-and-kernel-bypass:pmd}
A \emph{poll-mode driver}\index{poll-mode driver} runs in the application's process: it maps a card's rings into user space and
reads and refills them in a loop on a dedicated core, with no interrupts and no system calls. A \emph{user-space network
stack}\index{user-space network stack} implements IP, UDP and TCP in the application's process on top of such access to the card,
usually behind the standard socket calls, so that an unmodified program runs over it.
\end{definition}

\begin{remark}[The families of bypass]\label{rem:nw:network-cards-and-kernel-bypass:families}
Three kinds are in use, and a firm often runs two. Raw ring access through a vendor's interface (AMD's \texttt{ef\_vi} for its
Solarflare cards) or a portable poll-mode framework (DPDK) gives the program frames and nothing else: it must parse its own UDP and,
for orders, bring its own TCP. A user-space stack behind the socket calls (OpenOnload, open source, which intercepts the socket calls
of an unmodified binary) keeps the program's code and removes the kernel from its path. And Linux's own AF\_XDP moves frames from a
card's ring to a region of the program's memory through an in-kernel filter, in copy or in zero-copy mode, with the kernel still in
charge of the device. \Cref{met:nw:network-cards-and-kernel-bypass:choose} says which to use where.
\end{remark}

\begin{method}[Choosing a receive path]\label{met:nw:network-cards-and-kernel-bypass:choose}
\begin{enumerate}
\item Market data on the hot path: raw ring access, polled by a pinned core on the card's socket, the handler parsing UDP itself; the
  feed is multicast UDP, which needs no protocol state.
\item Order entry: a user-space TCP stack, or the card vendor's TCP in hardware (chapter~7), because order entry needs TCP's state and
  retransmission and writing one's own TCP is a project of its own.
\item Everything else (drop copy, reference data, monitoring): the kernel, on cores and rings of its own, so that it never shares a
  queue with the hot path.
\item In every case measure: the gain depends on the card, the host's tuning and the traffic, and a poller that is ever descheduled
  loses more than it gained.
\end{enumerate}
\end{method}

\omcode{code/firm/nicring/cpp/firm_nicring.hpp}{18}{37}{The receive ring of \texttt{firm.nicring}: the card fills descriptors it owns;
the poll loop processes those it owns and hands them back in batches.}\label{lst:nw:network-cards-and-kernel-bypass:ring}

\section{Zero copy and the transmit side}

\begin{definition}[Zero-copy networking]\label{def:nw:network-cards-and-kernel-bypass:zerocopy}
\emph{Zero-copy networking}\index{zero-copy networking} is the delivery of a packet to the program, or from it, in the buffer the card
reads or writes by direct memory access, without the processor copying its bytes between a kernel buffer and the program's.
\end{definition}

A copy of a 100-byte datagram is a few nanoseconds of memory traffic; what zero copy saves on the receive side is less the copy than
the machinery around it (the socket buffer, its accounting, the cache lines of a second location). On the transmit side the same
ring structure runs in reverse: the program writes an order into a transmit buffer, fills a descriptor and writes the card's doorbell
register, and the card fetches the frame by direct memory access. Each of those steps is a crossing of the processor's interconnect to
the card (One Quant Book~13, chapter~4). Some cards let the program push the frame's bytes into the card's own memory with processor
writes instead: AMD's \texttt{ef\_vi} documents such a programmed-I/O region, which ``avoids the latency associated with a DMA read'',
can be filled before the critical path and updated on it, and, in a cut-through mode, starts transmitting while the frame is still
crossing the bus. The cheapest transmit is the one prepared before it is needed: the order's
frame built in advance, with only price, size and identifier to fill, the same idea chapter~7 takes into hardware.

\section{Timestamps on the card}

The kernel can timestamp a packet when it enters the kernel (\texttt{SO\_TIMESTAMPING}, software receive timestamps), and a card with
its own clock can timestamp it when it reaches the card's port, a hardware timestamp in the sense of One Quant Book~13, chapter~5. The
difference is the time the packet spent in the card and its ring, which is exactly what a firm tuning its receive path wants to see and
what software timestamps cannot show. The card's clock is only as good as its synchronisation to the firm's time reference, the
subject of chapter~4, and the timestamps of the whole cage are compared in chapter~5.

\section{Tutorial: the kernel's receive path, measured and modelled}\label{tut:nw:network-cards-and-kernel-bypass:tut}

\begin{tutorial}
\textbf{Goal.} Measure what the kernel's receive path costs on the laptop, model what interrupts and polling cost on a card, and run
the \omterm{def:nw:network-cards-and-kernel-bypass:ring}{descriptor ring} in three languages. \textbf{End state:} \cref{tab:nw:network-cards-and-kernel-bypass:measured},
\cref{fig:nw:network-cards-and-kernel-bypass:measured,fig:nw:network-cards-and-kernel-bypass:coalesce}, and the ring's fixture green in
Python, C++ and Rust.
\begin{tutsteps}
\item \textbf{Three ways to read.} \texttt{cpp/nw\_rx.hpp} sends 64-byte datagrams over loopback, each carrying its send time, and reads
  them by a blocking \texttt{recvmsg}, by non-blocking reads in a loop, and by \texttt{recvmmsg} in batches of 16, with the kernel's software
  receive timestamp in each (\cref{lst:nw:network-cards-and-kernel-bypass:rx}). \texttt{python bench\_rx.py} runs 20\,000 datagrams per
  mode, 20 microseconds apart, and writes the quantiles.
\item \textbf{Split the path.} The kernel's timestamp splits each datagram's latency into getting into the kernel and getting out of it
  to the program.
\item \textbf{Model the card.} \texttt{nw\_nic.coalesce} and \texttt{poll} simulate a card's interrupts with coalescing and a polling
  core, with the model's costs; \texttt{python fig\_nic.py} writes \cref{fig:nw:network-cards-and-kernel-bypass:coalesce}'s data.
\item \textbf{The ring.} \texttt{firm.nicring} runs the same stream of arrivals and polls through four ring configurations in Python,
  C++ and Rust and checks counters and hashes against the shared fixture.
\end{tutsteps}
\textbf{What to change next.} Pin the busy-polling reader to a core and the sender to another (\texttt{taskset}) and read its tail
again; in the ring's fixture, find the smallest ring that drops nothing when the host hands back in batches of 32.
\end{tutorial}

\omcode{code/networks/03-network-cards-and-kernel-bypass/cpp/nw_rx.hpp}{100}{116}{The three reads: a blocking read, a non-blocking
read in a loop, and a batched read; each datagram keeps its send, kernel and application times.}\label{lst:nw:network-cards-and-kernel-bypass:rx}

\begin{table}[ht]
\centering\small
\begin{tabular}{lrrrr}
\toprule
Read method & Into the kernel, p50 & Kernel to program, p50 & Total, p50 & Total, p99 \\
\midrule
Blocking \texttt{recvmsg} & \qty{0.75}{\micro\second} & \qty{15.4}{\micro\second} & \qty{16.2}{\micro\second} & \qty{215}{\micro\second} \\
Busy polling & \qty{0.62}{\micro\second} & \qty{0.92}{\micro\second} & \qty{1.5}{\micro\second} & \qty{3\,457}{\micro\second} \\
Batched \texttt{recvmmsg} & \qty{0.45}{\micro\second} & \qty{21.4}{\micro\second} & \qty{22.0}{\micro\second} & \qty{75}{\micro\second} \\
\bottomrule
\end{tabular}
\caption{The kernel's receive path over loopback, 20\,000 datagrams of 64 bytes per method, 20 microseconds apart (the batched
reader's sender sends bursts of 16). Measured on a laptop (Intel Core Ultra 7 155H) under WSL2, no isolated cores; g++~11,
\texttt{-O2}. Data: \texttt{bench\_rx.py}, \texttt{measured\_rx.csv}.}\label{tab:nw:network-cards-and-kernel-bypass:measured}
\end{table}

\begin{omfigure}
\begin{tikzpicture}
\begin{axis}[width=11.5cm,height=5.4cm,ymode=log,xmin=0,xmax=1,ymin=0.7,ymax=6000,
  xlabel={quantile},ylabel={latency (\unit{\micro\second}, log scale)},
  ytick={1,10,100,1000},yticklabels={1,10,100,1\,000},xtick={0,0.25,0.5,0.75,0.9,1},
  xticklabels={0,0.25,0.5,0.75,0.9,1},
  tick label style={font=\small},label style={font=\small},grid=major,grid style={gray!25},
  legend columns=3,legend style={font=\footnotesize,at={(0.5,-0.3)},anchor=north,draw=none,fill=none,
  /tikz/every even column/.append style={column sep=8pt}},table/col sep=comma]
\addplot[omThm,very thick,mark=*,mark size=1.3pt] table[x=q,y=blocking]{figdata/networks/03-network-cards-and-kernel-bypass/measured_rx_q.csv};
\addplot[omDef,very thick,mark=square*,mark size=1.3pt] table[x=q,y=busy]{figdata/networks/03-network-cards-and-kernel-bypass/measured_rx_q.csv};
\addplot[omMeth,very thick,mark=triangle*,mark size=1.6pt] table[x=q,y=batched]{figdata/networks/03-network-cards-and-kernel-bypass/measured_rx_q.csv};
\legend{blocking read,busy polling,batched read}
\end{axis}
\end{tikzpicture}

\omcaption{Quantiles of the time from \texttt{sendto} to the program holding the datagram, on the laptop's loopback, by read method.
Busy polling cuts the median tenfold; without an isolated core the polling thread is sometimes descheduled for milliseconds, and every
datagram behind it waits: from the 95th quantile on, it is the worst of the three. Measured on a laptop (Intel Core Ultra 7 155H) under
WSL2, no isolated cores. Data: \texttt{bench\_rx.py}.}\label{fig:nw:network-cards-and-kernel-bypass:measured}
\end{omfigure}

The measurement shows the two halves of this chapter's argument (\cref{tab:nw:network-cards-and-kernel-bypass:measured}). Getting a
datagram into the kernel costs under a microsecond; getting it out to a sleeping program costs fifteen at the median and hundreds of
microseconds at the 99th percentile, because the program must be woken and scheduled. A program that never sleeps removes that: busy
polling brings the median down to \qty{1.5}{\micro\second}. And a program that never sleeps must never be interrupted: on a laptop
whose cores are shared with everything else the poller is sometimes descheduled, and the datagrams that arrive meanwhile wait for
milliseconds. Batching amortises the system call but not the wake-up, and waits for its batch. Loopback involves no card at all, so
these figures are the kernel's share alone; a real card adds its ring, its interrupt and its coalescing.

\omcode{code/networks/03-network-cards-and-kernel-bypass/python/nw_nic.py}{24}{45}{The coalescing model: an interrupt at the $m$-th
waiting packet or $u$ after the first, then wake-up, interrupt cost and per-packet cost.}\label{lst:nw:network-cards-and-kernel-bypass:model}

\begin{omfigure}
\begin{tikzpicture}
\begin{groupplot}[group style={group size=2 by 1,horizontal sep=1.8cm},width=6.6cm,height=5.0cm,xmin=0,xmax=66,
  xlabel={coalescing timer (\unit{\micro\second})},tick label style={font=\small},label style={font=\small},
  grid=major,grid style={gray!25},table/col sep=comma,unbounded coords=jump]
\nextgroupplot[ymin=0,ymax=48,ylabel={mean latency (\unit{\micro\second})},xtick={0,8,16,32,64},
  legend columns=2,legend style={font=\footnotesize,at={(1.15,-0.32)},anchor=north,draw=none,fill=none,
  /tikz/every even column/.append style={column sep=8pt}}]
\addplot[omDef,very thick,mark=*,mark size=1.3pt] table[x=timer_us,y=lat_01]{figdata/networks/03-network-cards-and-kernel-bypass/coalesce.csv};
\addplot[omThm,very thick,mark=square*,mark size=1.3pt] table[x=timer_us,y=lat_05]{figdata/networks/03-network-cards-and-kernel-bypass/coalesce.csv};
\legend{100\,000 packets a second,500\,000 packets a second}
\nextgroupplot[ymin=0,ymax=0.7,ylabel={share of one core},xtick={0,8,16,32,64}]
\addplot[omDef,very thick,mark=*,mark size=1.3pt] table[x=timer_us,y=cpu_01]{figdata/networks/03-network-cards-and-kernel-bypass/coalesce.csv};
\addplot[omThm,very thick,mark=square*,mark size=1.3pt] table[x=timer_us,y=cpu_05]{figdata/networks/03-network-cards-and-kernel-bypass/coalesce.csv};
\end{groupplot}
\end{tikzpicture}

\omcaption{Model: \omterm{def:nw:network-cards-and-kernel-bypass:coalesce}{interrupt coalescing} at two packet rates. Longer timers cost latency and save CPU. At 500\,000 packets a second an
interrupt per packet saturates the core (no point at zero). A polling core, not drawn, holds latency near
\qty{0.4}{\micro\second} at both rates and uses the whole core. Model costs: \qty{2}{\micro\second} per interrupt,
\qty{3}{\micro\second} of wake-up, \qty{0.3}{\micro\second} per packet. Data: \texttt{fig\_nic.py}.}\label{fig:nw:network-cards-and-kernel-bypass:coalesce}
\end{omfigure}

The model (\cref{lst:nw:network-cards-and-kernel-bypass:model}) takes \qty{2}{\micro\second} of CPU per interrupt,
\qty{3}{\micro\second} of wake-up and \qty{0.3}{\micro\second} per packet: model values of the order of the laptop's wake-up, not a
particular card's. At 100\,000 packets a second an interrupt per packet gives \qty{5.6}{\micro\second} of mean latency for 23\% of a
core; an \qty{8}{\micro\second} timer gives \qty{11.7}{\micro\second} for 14\%; polling gives \qty{0.4}{\micro\second} for the whole
core. At 500\,000 an interrupt per packet is beyond one core (the saturation rate of
\cref{prop:nw:network-cards-and-kernel-bypass:coalesce} is 435\,000 a second), and coalescing is what keeps the host alive. A trading
host resolves the trade-off by not choosing it: the hot path polls on cores bought for the purpose, and everything else coalesces.

\section{Build: the descriptor ring}\label{bld:nw:network-cards-and-kernel-bypass:nicring}

\begin{build}
\bfield{Purpose} A model of a card's receive ring precise enough to reason about drops, batching and polling budgets, in the three
languages of the firm's systems.
\bfield{Interface} Python \texttt{firm\_nicring.Ring(size, refill)} with \texttt{nic\_rx(seq, length)}, \texttt{poll(budget)},
\texttt{run(events)} and the counters \texttt{rx}, \texttt{drops}, \texttt{processed}, \texttt{doorbells}, \texttt{max\_owned},
\texttt{hash}; \texttt{make\_events}. C++20 \texttt{firm::nicring::Ring} (\texttt{cpp/firm\_nicring.hpp}) and Rust
\texttt{firm\_nicring::Ring} with the same methods.
\bfield{Rules} The ring's size is a power of two; the card never overwrites a descriptor the host owns; the host processes in order,
never more than its budget per poll, and hands back in batches of \texttt{refill}, one doorbell per batch; the hash is FNV-1a over
the processed packets' sequence numbers and lengths in processing order.
\bfield{Acceptance tests} \texttt{code/firm/nicring/}: the shared fixture (20\,000 arrivals, four configurations) reproduced exactly in
Python, C++ and Rust; a full ring drops and a processed but unreturned descriptor still counts as full; under random interleavings
every accepted packet is processed once and in order.
\bfield{Stretch} A transmit ring with a completion queue; several rings with a steering function, to reproduce the imbalance of a hash
that puts two busy groups on one ring.
\end{build}

\begin{omsources}
\item Linux kernel documentation: ``NAPI'', ``Scaling in the Linux Networking Stack'', ``Timestamping'', ``AF\_XDP'';
\texttt{include/uapi/linux/ethtool.h} (\texttt{struct ethtool\_coalesce}); man page \texttt{recvmmsg(2)}.
\item DPDK Programmer's Guide, ``Poll Mode Driver''; AMD, \emph{ef\_vi User Guide}; OpenOnload README.
\end{omsources}

\section{Exercises}

\begin{exercise}[$\star$]\label{exo:nw:network-cards-and-kernel-bypass:1}
A ring of 512 descriptors, of which the host holds 40, receives 1.5 million packets a second while its core is interrupted for
\qty{400}{\micro\second}. How many packets are dropped?
\end{exercise}

\begin{exercise}[$\star$]\label{exo:nw:network-cards-and-kernel-bypass:2}
With coalescing at 20 microseconds, and packets at 200\,000 a second, how many packets does an interrupt serve on average, and at what
rate do interrupts arrive?
\end{exercise}

\begin{exercise}[$\star$]\label{exo:nw:network-cards-and-kernel-bypass:3}
From \cref{tab:nw:network-cards-and-kernel-bypass:measured}, what share of the blocking read's median is spent after the datagram entered
the kernel?
\end{exercise}

\begin{exercise}[$\star\star$]\label{exo:nw:network-cards-and-kernel-bypass:4}
Why does \omterm{def:nw:network-cards-and-kernel-bypass:rss}{receive-side scaling}'s even spreading of flows work against a feed handler, and what does \omterm{def:nw:network-cards-and-kernel-bypass:rss}{flow steering} change?
\end{exercise}

\begin{exercise}[$\star\star$]\label{exo:nw:network-cards-and-kernel-bypass:5}
With the model's costs, at what packet rate does an interrupt per packet use a whole core? And if every interrupt serves 16 packets?
\end{exercise}

\begin{exercise}[$\star\star$]\label{exo:nw:network-cards-and-kernel-bypass:6}
Busy polling improved the median tenfold and made the 99th percentile sixteen times worse
(\cref{tab:nw:network-cards-and-kernel-bypass:measured}). Explain both, and say what must change on a production host.
\end{exercise}

\begin{exercise}[$\star\star\star$]\label{exo:nw:network-cards-and-kernel-bypass:7}
\emph{Coding.} With \texttt{firm.nicring} and the fixture's events, find the smallest ring size that drops nothing when descriptors are
handed back in batches of 32.
\end{exercise}

\begin{exercise}[$\star\star\star$]\label{exo:nw:network-cards-and-kernel-bypass:8}
\emph{Find the flaw.} ``We moved order entry to a raw poll-mode interface and wrote our own minimal TCP: no retransmission timer, since
colocation links do not lose packets. Latency fell by two microseconds.''
\end{exercise}

\section{Problem: Interrupts or Polling}

\begin{problem}[{Weekend problem --- the price of a sleeping reader}]\label{pb:nw:network-cards-and-kernel-bypass:1}
A feed arrives at 200\,000 packets a second on one receive ring. Use the chapter's model costs: \qty{2}{\micro\second} of CPU per
interrupt, \qty{3}{\micro\second} of wake-up, \qty{0.3}{\micro\second} per packet, and a polling loop that checks the ring every
\qty{0.2}{\micro\second}.

\medskip\noindent\textbf{Part I --- One interrupt per packet.}
\begin{enumerate}
\item What share of a core do interrupts and packets take?
\item What is the smallest latency a packet can have?
\item The model's mean latency is \qty{6.3}{\micro\second}. Why is it above the minimum?
\item At what rate would the core saturate?
\end{enumerate}

\medskip\noindent\textbf{Part II --- A 10-microsecond timer.}
\begin{enumerate}[resume]
\item How many packets does an interrupt serve on average?
\item At what rate do interrupts come?
\item What share of the core is used?
\item The model's mean latency is \qty{12.4}{\micro\second}. Where does the increase come from?
\end{enumerate}

\medskip\noindent\textbf{Part III --- Polling.}
\begin{enumerate}[resume]
\item What is the mean wait for the next check of the ring?
\item What is the mean latency, and the model's (\qty{0.41}{\micro\second})?
\item What share of a core is used?
\item What must be true of that core?
\end{enumerate}

\medskip\noindent\textbf{Part IV --- The verdict.}
\begin{enumerate}[resume]
\item State the \emph{named result}: the mean latency and CPU of the timer and of polling at this rate, and the rate at which interrupts
  alone would take a whole core.
\item How many microseconds per packet does polling save over the timer?
\item At what packet rate does the timer's CPU reach a quarter of a core?
\item Why do trading hosts poll the hot path and coalesce the rest?
\item What would you measure on the real card before believing the model?
\item What does a \omterm{def:nw:network-cards-and-kernel-bypass:ring}{descriptor ring} of 1\,024 entries give the polling core, in time, at this rate?
\item In one sentence: what does a sleeping reader cost?
\item Which of the chapter's three families of bypass would you use for this feed, and why?
\end{enumerate}
\end{problem}

\section{Interview questions}

\begin{interviewq}[$\star$ \iqroles{developer}]\label{iq:nw:network-cards-and-kernel-bypass:1}
Walk me through what happens between a packet arriving at the card and \texttt{recv} returning.
\end{interviewq}

\begin{interviewq}[$\star\star$ \iqroles{developer}]\label{iq:nw:network-cards-and-kernel-bypass:2}
What is \omterm{def:nw:network-cards-and-kernel-bypass:coalesce}{interrupt coalescing}? How would you set it on a market-data host?
\end{interviewq}

\begin{interviewq}[$\star\star$ \iqroles{developer}]\label{iq:nw:network-cards-and-kernel-bypass:3}
What is kernel bypass, what forms does it take, and what do you give up?
\end{interviewq}

\begin{interviewq}[$\star\star$ \iqroles{developer}]\label{iq:nw:network-cards-and-kernel-bypass:4}
Your busy-polling feed handler has a great median and a terrible 99.9th percentile. What do you check?
\end{interviewq}

\begin{interviewq}[$\star\star\star$ \iqroles{developer}]\label{iq:nw:network-cards-and-kernel-bypass:5}
A \omterm{def:nw:network-cards-and-kernel-bypass:ring}{descriptor ring} drops packets although the handler's average load is 30\%. Explain how, and how you would size the ring.
\end{interviewq}

\begin{interviewq}[$\star\star$ \iqroles{developer, researcher}]\label{iq:nw:network-cards-and-kernel-bypass:6}
Why would you want hardware receive timestamps when you already timestamp in the kernel?
\end{interviewq}
